%Matz JB 2012
%Using Brittish spelling of realise, colour, and conjunctions


%%%%%%%%%%%%%%%%%%%%%%%%%%%%%%%%%%%%%%%%%%%%%%%%%%%%%%%%%%%%%%%%%%%%%%

%a4paper

\PassOptionsToPackage{table}{xcolor}
\documentclass[a4paper, 11pt]{article}
\usepackage[T1]{fontenc}
\usepackage{tocloft}
\usepackage{fix-cm}
\usepackage{titlesec}

\usepackage{tikz} % Style 2, 3, 4
\usetikzlibrary{calc} % 2, 3, 4
\usetikzlibrary{decorations,decorations.shapes,shapes} % 3
\usepackage{float}
\usepackage[top=2cm, bottom=2cm]{geometry}
\usepackage[bookmarks]{hyperref}
\usepackage{geometry}
\usepackage{fancyhdr}
\usepackage{tikz}
\usepackage{amsmath}%för bmatrix
\usepackage{url}
\usepackage{eufrak}
\usepackage{color, colortbl}
\usepackage{pdfpages}
\usepackage{graphicx}
\usepackage{listings}
\usepackage{xkeyval}
\usepackage{xfor}
\usepackage{etoolbox}
\usepackage{longtable}
\usepackage{supertabular}
\usepackage{ textcomp }
\definecolor{lightblue}{rgb}{0.9,0.95,0.99}
\usepackage{graphics}
\usepackage[sanitize=none, nonumberlist, acronym, toc, translate=false]{glossaries}
\usepackage{glossary-superragged}
\glsdisablehyper %no links
\usepackage{tikz} % Style 2, 3, 4
\usetikzlibrary{calc} % 2, 3, 4
\usetikzlibrary{decorations,decorations.shapes,shapes} % 3

\usepackage{tikz}
\newcommand*\circled[1]{\tikz[baseline=(char.base)]{%base
            \node[shape=circle,draw,inner sep=0.9pt] (char) {\footnotesize\bf \texttt{#1}};}}



%http://tex.stackexchange.com/questions/2275/keeping-tables-figures-close-to-where-they-are-mentioned

\usepackage{afterpage} %10i8~12

%http://www.latex-community.org/viewtopic.php?f=5&t=2423
    \usepackage[font=small, labelfont={color=labelcolor, bf}]{caption}
    \captionsetup{%
      within=section
}


%\restylefloat{figure}%place figure here!

%listing tabs
\def \tabsize {5}
\def \basicstyle {\footnotesize} % {\scriptsize} %small footnotesize



%\afterpage{\clearpage} %interesting

% See p.105 of "TeX Unbound" for suggested values.
    % See pp. 199-200 of Lamport's "LaTeX" book for details.
    %   General parameters, for ALL pages:
    \renewcommand{\topfraction}{0.95}	% max fraction of floats at top
    \renewcommand{\bottomfraction}{0.2}	% max fraction of floats at bottom
    %   Parameters for TEXT pages (not float pages):
    \setcounter{topnumber}{2}
    \setcounter{bottomnumber}{0}
    \setcounter{totalnumber}{2}     % 2 may work better
%    \setcounter{dbltopnumber}{2}    % for 2-column pages
    \renewcommand{\dbltopfraction}{0.9}	% fit big float above 2-col. text
    \renewcommand{\textfraction}{0.17}	% allow minimal text w. figs
    %   Parameters for FLOAT pages (not text pages):
    \renewcommand{\floatpagefraction}{0.7}	% require fuller float pages
	% N.B.: floatpagefraction MUST be less than topfraction !!
    \renewcommand{\dblfloatpagefraction}{0.7}	% require fuller float pages



%\usepackage{refcheck} %To find references (figures) I have not referenced

%%%%%%%%%%%%%%%%%%%%%%%%%%%%%%%%%%%%%%%%%%%%%%%%%%%%%%%%%%%%%%%%%%%%%%

%http://cs.gmu.edu/~sean/stuff/DoingLatexRight/
\usepackage{times} %30/7 ~12

\usepackage[
    %backend = biber,%fungerade inte ihop med texlua. texlua "hittar inte bib-filen" 21/5
%style=numeric,
%    style = authoryear-icomp,
    style = alphabetic, 
%style=superscript,   
    %sortlocale = de_DE,
    natbib = true,
    url = true, 
    doi = true,
    %sorting = sortinit,
    eprint = true % vad gör denna?
]{biblatex}

%\usepackage[super]{natbib}
%\addbibresource{refs.bib}


\usepackage{todonotes}
%\usepackage[colorinlistoftodos, textwidth=2cm,% backgroundcolor=light_blue,
%            textsize=scriptsize,disable]{todonotes} %if lines start to go vertical, compile again... [25i10~09]
%http://midtiby.blogspot.com/2007/09/todo-notes-in-latex.html

\definecolor{shade}{HTML}{CCECCD}%must be after the "todonotes" package
\definecolor{shade_q}{HTML}{D4D5E5}%must be after the "todonotes" package

\newcommand{\todonum}[1] %to get rid of the lines
{\todonumm{#1}}

%for todo numberings
\newcounter{todocounter}
\newcommand{\todonumm}[2][]
{\stepcounter{todocounter}\todo[inline,#1]{\thetodocounter: #2}}

\newcommand{\subsubsubsection}[1]
{\paragraph{#1}}



\setcounter{tocdepth}{3}


\usepackage[section]{algorithm}
\usepackage{algorithmicx}
\usepackage{algpseudocode} %[noend]
\usepackage{multicol}
\usepackage{pdfcomment}
\usepackage{epsfig}
\usepackage{epic} %this may or may not be able to download from servers (12 sept ~09), yes now it works [13i9~09]
\usepackage{eepic}

%\usepackage[mediumspace,mediumqspace,squaren,binary]{SIunits}
\usepackage[squaren,binary]{SIunits}


\usepackage{graphicx}
%\usepackage[table]{xcolor}% [26i10~09] http://www.cv-templates.info/2009/07/alternate-row-shading-latex/
\usepackage{nicefrac}%funkar inte för Ubuntu och jag hittar inte skiten 28i9~11

\usepackage{wrapfig}
%\usepackage{mathptmx}

%For line numbers in the margin
%\usepackage{lineno,color}
%\linenumbers

%http://www.terminally-incoherent.com/blog/2007/09/19/latex-squeezing-the-vertical-white-space/
\renewenvironment{itemize}%
  {\begin{list}{\arabic{enumi}.}%
     {\topsep=0.1in\itemsep=0.5Em\parsep=0in\partopsep=0in\usecounter{enumi}}%
   }{\end{list}}

%\newcommand{\R}{\ensuremath{\mathcal{R}}}

\newcommand{\rd}{\ensuremath{\mathrm{d}}}
\newcommand{\id}{\ensuremath{\,\rd}}

%Making symbols bold in math environment
\newcommand{\mt}[1]{\boldsymbol{\mathbf{#1}}}

\usepackage{verbatim}
\usepackage[english, swedish]{babel}
\usepackage{blindtext}
\usepackage{microtype}

\definecolor{keywords}{RGB}{255, 0, 90}
\definecolor{comments}{RGB}{60, 179, 113}
\definecolor{tableShade}{HTML}{E0FFE0}   %iTunes
\definecolor{tableShade2}{HTML}{E0FFE0} %Finder


\definecolor{keywordss}{rgb}{0, .4, 0}
\definecolor{commentss}{rgb}{.3, .3, .7}
\definecolor{stringss}{rgb}{.3, .7, .3}

\definecolor{keycolor}{rgb}{0, 0, 0.4}

\definecolor{light_blue}{rgb}{.8, .8, 1}
\definecolor{blue}{rgb}{0, 0, .6}
\definecolor{green}{rgb}{0, .5, 0}
\definecolor{red}{rgb}{.7, 0, 0}
\definecolor{LightSkyBlue}{rgb}{0.6, 0.6, 0.9}
\definecolor{MidnightBlue}{rgb}{0.2, 0.2, 0.9}
\definecolor{darkblue}{rgb}{0.8, 0, 0.5} %test

\definecolor{labelcolor}{rgb}{0.1, 0.3, 0.1}

%\definecolor{gray}{rgb}{0.9, 0.95, 1}
\definecolor{gray}{rgb}{0.5, 0.5, 0.5}


\usepackage{abstract}
\usepackage{bm}


\usepackage{textcomp}%http://tex.stackexchange.com/questions/55036/package-listings-error-undefined-control-sequence-because-of-single-quote-sy

%\usepackage{underscore}
\usepackage[strings]{underscore}

\setlength{\parindent}{0Em}%4 letters indentation (ex=letter?!)
\setlength{\bibnamesep}{0.5Em} %bibliography separation of different bibliography items


\usepackage[parfill]{parskip}
\hyphenpenalty=500
%\setlength{\parskip}{0cm}

%http://robjhyndman.com/researchtips/squeezing-space-with-latex/

%\usepackage[compact]{titlesec}

\usepackage{graphics}

%\usepackage[explicit]{titlesec}
\usepackage{titlesec}
\begin{comment}
\titleformat{\section}[display]
{\bfseriese}
{\hfill \tikz[remember picture] \node[] (nr) {\fontsize{120}{70}\selectfont\color{green!50!white}\textbf{\thesection}};
\begin{tikzpicture}[overlay,remember picture]
\coordinate (leftborder) at ($(nr)-(100,0)$);
\coordinate (left) at ($(nr.west)-(1.5,0)$);
\draw[line width=0.5em, draw=green!50!white] ($(nr.north east)+(0,0.4)$) -- ($(nr.north west)+(0,0.4)$) -- (left) -- (leftborder);
\end{tikzpicture}}
{-2ex}
{\filleft\fontsize{40}{70}\selectfont\scshape}
[\vspace{0ex}]
\end{comment}

\begin{comment}
\titleformat{\section}[display]
{\bfseriese}
{
\hfill \tikz[remember picture] \node[] (nr) {\fontsize{100}{40}\selectfont\color{black}\thesection};
\begin{tikzpicture}[overlay, remember picture]
\draw[line width = 2em, draw=black] (0, 1.5) -- (10, 1.5);
\end{tikzpicture}}
{-1ex}
{\filright\fontsize{20}{20}\selectfont\scshape}
[\vspace{0ex}]
\end{comment}


\newcommand*{\justifyheading}{\raggedleft}


\begin{comment}
%right adjusted:
\titleformat{\chapter}[display]{\normalfont\huge\bfseries\justifyheading}{\chaptertitlename\ \thechapter}{20pt}{\Huge}

\titleformat{\section}{\normalfont\Large\bfseries\raggedleft}{\thesection}{1em}{}

\titleformat{\subsection}{\normalfont\large\bfseries\justifyheading}{\thesubsection}{1em}{}
\end{comment}

%\newcommand*{\justifyheading}{\raggedleft}


\begin{comment}
\titleformat{\section}
{\titlerule
\vspace{.8ex}%
\normalfont\itshape}
{\thesection.}{.5em}{}
\end{comment}

%\usepackage[linedheaders,parts,pdfspacing]{classicthesis} % ,manychapters


%\draw[decoration={shape backgrounds,shape size=.5cm,shape=signal},signal from=west, signal to=east,decorate, draw=red!50!black, fill=red!50, decoration={shape sep=.5cm},line 

%\titlespacing{\section}{0pt}{2ex}{1ex}
%\titlespacing{\subsection}{0pt}{1ex}{0ex}
%\titlespacing{\subsubsection}{0pt}{0.5ex}{0ex}

\begin{comment}
\titleformat{\section}[display]
    {\normalfont\Large\raggedleft}
    {\MakeUppercase{\sectiontitlename}%
    \rlap{ \resizebox{!}{1.2cm}{\thesection} \rule{15cm}{1.2cm} } }
    {10pt}{\Huge}

\titlespacing*{\section}{0pt}{30pt}{20pt}
\end{comment}



%\titleformat{name=\section,page=even}[leftmargin]

% {\filleft\scshape}{\thesection}{.5em}{}

\lstset{linewidth=0.465\textwidth}%width of code box, 4in
\lstset{xleftmargin=0.1in}
\lstset{backgroundcolor=\color{white}}

%keywordstyle causes acrobat reader to fuck up

\def \lstsetC {
\lstset{
	language = C++,
frame=bottomline,
	keywordstyle =\bfseries\ttfamily\color[rgb]{0, 0, 0.4},
	identifierstyle =\bfseries\ttfamily\color[rgb]{0, 0.4, 0},
	commentstyle =\color[rgb]{0.333,0.345,0.333},
	stringstyle =\ttfamily\color[rgb]{0.627,0.126,0.941},
	showstringspaces = false,
	basicstyle = \basicstyle, %\small,
	tabsize = \tabsize,
	prebreak = \raisebox{0ex}[0ex][0ex]{\ensuremath{$\ldots$}},
	breakatwhitespace =false,
	aboveskip ={1.5\baselineskip},
        upquote =true,
        captionpos =t,
        extendedchars =true,
morekeywords={omp, parallel, \#pragma, for, reduction, pragma, default, QueryPerformanceCounter, QueryPerformanceTimer, QueryPerformanceFrequency, LARGE_INTEGER}
}
}

%%error below

\def \lstsetMATLAB {
\lstset{
	language = Matlab,
	keywordstyle =\bfseries\ttfamily\color[rgb]{0, 0, 0.4},
	identifierstyle =\bfseries\ttfamily\color[rgb]{0,0.4,0},
	commentstyle =\color[rgb]{0.333, 0.333, 0.333},
	stringstyle =\ttfamily\color[rgb]{0.1, 0.2, 0.1},
	showstringspaces = false,
	basicstyle = \basicstyle,
%	numberstyle=\footnotesize,
        %numberstyle       = \footnotesize\color{darkbrown}\slshape,
%	numbers=left,
%	stepnumber =1,
	%numbersep =10pt,
	tabsize    = 5,
%       frame = l,
        breaklines = true,
	prebreak   = \raisebox{0ex}[0ex][0ex]{\ensuremath{$\ldots$}},
	breakatwhitespace =false,
	aboveskip     = {1.5\baselineskip},
        %columns =fullflexible,
        upquote       = true,
        captionpos    = t,
        extendedchars = true,
        morekeywords={try, catch, rethrow, complex, tic, toc},
        morecomment=[l][keywordstyle]{!}
}
}



\def \lstsetCUDA {
\lstset{
	language         = C,
	keywordstyle     = \bfseries\ttfamily\color[rgb]{0, 0, 0.4},
	identifierstyle  = \bfseries\ttfamily\color[rgb]{0, 0.4, 0},
	commentstyle     = \color[rgb]{0.133, 0.3, 0.133},
	stringstyle      = \ttfamily\color[rgb]{0.1, 0.2, 0.1},
	showstringspaces = false,
	basicstyle       = \basicstyle,
%	numberstyle=\footnotesize,
        %numberstyle       = \footnotesize\color{darkbrown}\slshape,
%	numbers=left,
%	stepnumber =1,
	%numbersep =10pt,
%        backgroundcolor=\color[rgb]{0.95, 0.95, 0.95},
	tabsize       = 1,
        rulecolor=\color{black}, 
%        frame = l,
        breaklines    = true,
	prebreak      = \raisebox{0ex}[0ex][0ex]{\ensuremath{$\ldots$}},
	breakatwhitespace =false,
	aboveskip     = {1.5\baselineskip},
        %columns =fullflexible,
        upquote       = true,
        captionpos    = t,
        extendedchars = true,
        morekeywords={QueryPerformanceCounter, dim3, __global__, cudaEvent_t, cudaEventRecord, cudaEventStop, cudaEventSynchronize,cudaEventElapsedTime, cudaEventCreate, cudaThreadSynchronize},
morecomment=[s]{<<<}{>>>}
}
}

%\lstset{frame=lrb,xleftmargin=\fboxsep,xrightmargin=-\fboxsep}



\newcommand{\code}[1]
{\texttt{\textbf{#1}}}


%updated 22i11~11 latest version of all
%this version will work with twocolumns
\newcommand{\addfigure}[4]
{
\begin{figure}[tbh]

  \begin{center}
        \includegraphics[width=#2]{#1}
  \end{center}

  \caption{#3}
\label{fig:#4}
  \vspace{-5pt}
\end{figure}

}

%updated 22i11~11 latest version of all
%this version will work with twocolumns
%updated with makebox, to make the figures centered without regard of textwidth
\newcommand{\addfig}[5]
{
\vspace{1Em}
\begin{figure}[tbh]
  \makebox[\textwidth][c]{\includegraphics[width=#2]{#1}}%
  \caption[#5]{#3}
  \label{fig:#4}
%\vspace{1Em}

\end{figure}
}




\begin{comment}

\begin{center}
\begin{figure}[!h]

%  \begin{center}
        \includegraphics[width = #2]{#1}
 % \end{center}

  \caption[#5]{#3}
\label{fig:#4}
  \vspace{-5pt}
\end{figure}
\end{center}
}
\end{comment}



%Matz Johansson B. 20i9~09, 14:30
\newcommand{\ikon}[1]
{
  \def\iheight{1em} %new variable containing the height of the icons based
  % on the height of 'M' in the current font [20i9~09]
  \raisebox{-0.2\iheight} %we lower the graphics under baseline so an icon [] will between...
  {                     %parenthesis look vertically centered
    \hspace{-1.7em}
    {
      \includegraphics[height=\iheight]{#1}
    } %hspace, no empty lines in renewcommand or it will introduce newlines
    \hspace{-1em}%av nån anledning måste man göra såhär [24i10~09]
  } %raisebox
} %newcommand


%Matz Johansson B. 20i9~09, 14:30
\newcommand{\ikonn}[1]
{
  \def\iheight{15pt} %new variable containing the height of the icons based
  % on the height of 'M' in the current font [20i9~09]
  \raisebox{-0.2\iheight} %we lower the graphics under baseline so an icon [] will between...
  {                     %parenthesis look vertically centered
    \hspace{-2em}
    {
	\vspace{0Em}
	{
      \includegraphics[height=\iheight]{#1}
    }
	} %hspace, no empty lines in renewcommand or it will introduce newlines
    \hspace{-1em}%av nån anledning måste man göra såhär [24i10~09]
  } %raisebox
} %newcommand




%\hypersetup{colorlinks=false, linkcolor=black}%pdfborder

\hypersetup{hidelinks=true} %change this...
%
\hypersetup{
  colorlinks,
  citecolor = softgreen,%citecol,
  linkcolor = linkcol, %black
  urlcolor  = whiteo}

\hypersetup{
  citebordercolor = white,
  filebordercolor = white,
  linkbordercolor = white
}

%\hypersetup{backref=true}

%vi vill inte se boxar runt länkar [19i8~11]
%\usepackage[colorlinks=true]{hyperref}

%\usepackage[colorlinks=true, linktocpage=true]

\newcommand{\HRule}{\rule{\linewidth}{0.5mm}}

%\begin{comment}
%från Christian Von Schultz' sida
\usepackage{bbm}
\newcommand{\N}{\ensuremath{\mathbbm{N}}}
\newcommand{\Z}{\ensuremath{\mathbbm{Z}}}
\newcommand{\Q}{\ensuremath{\mathbbm{Q}}}
\newcommand{\R}{\ensuremath{\mathbbm{R}}}
\newcommand{\C}{\ensuremath{\mathbbm{C}}}


%Matz transpose
\newcommand{\tr}{\top}
%\newcommand{\tr}{\mathrm T}


%\end{comment}


%27/6 1012
\newcommand{\com}[1]
{
\begin{comment} #1 \end{comment}
}


\newcommand{\kommentar}[1]
 {
\pdfmargincomment[open=true, icon=note, color=llightgray, author=Matz]{#1}
 }
% {  \footnotesize\colorbox{lightblue}{#1}}

%\renewcommand*{\glstextformat}[1]{\textsf{#1}}


\lstset{language    = C}
\lstset{linewidth   = 5.78in}%width of code box 4.5in
\lstset{xleftmargin = 0.2in}%0.25

\lstset{keywordstyle=\color{keywords}\bfseries}
\lstset{commentstyle=\color{DarkGreen}}
%\lstset{backgroundcolor=\color{CodeBackground}}

\definecolor{darkblue}{rgb}{0, 0, 0.35}
\definecolor{darkgreen}{rgb}{0, 0.1, 0}
\definecolor{darkbrown}{rgb}{0.6, 0.6, 0}

\definecolor{softgreen}{rgb}{0.8, 0, 0.8}
\definecolor{linkcol}{rgb}{0, 0, 0.7} %very visible to be able to see how many the report has over the report

\definecolor{darkgray}{rgb}{0.35, 0.35, 0.35}
\definecolor{pink}{rgb}{0.35, 0.35, 0.2}
\definecolor{whiteo}{rgb}{0.1, 0.1, 0.8}

\definecolor{white}{rgb}{1, 1, 1}
%\definecolor{citecol}{rgb}{0.2, 0.2, 1}
\definecolor{citecol}{rgb}{1, 0.2, 0.2} %very visible to be able to see how many the report has over the report



%\color[rgb]{0,0.5,1}

\lstset{
	language = Matlab,
	keywordstyle =\bfseries\ttfamily\color[rgb]{0, 0, 0.4},
	identifierstyle =\ttfamily,
	commentstyle =\color[rgb]{0.133,0.545,0.133},
	stringstyle =\ttfamily\color[rgb]{0.627,0.126,0.941},
	showstringspaces = false,
	basicstyle = \basicstyle,
%	numberstyle=\footnotesize,
        %numberstyle       = \footnotesize\color{darkbrown}\slshape,
%	numbers=left,
%	stepnumber =1,
	%numbersep =10pt,
	tabsize =\tabsize,
       % frame = l,
        breaklines =true,
	prebreak = \raisebox{0ex}[0ex][0ex]{\ensuremath{$\ldots$}},
	breakatwhitespace =false,
	aboveskip ={1.5\baselineskip},
        %columns =fullflexible,
        upquote =true,
        captionpos =t,
        extendedchars =true
}



%  \usepackage{caption}
%\DeclareCaptionFont{white}{\color{white}}
\DeclareCaptionFont{white}{\color{white}} %black
\DeclareCaptionFormat{listing}{\colorbox[rgb]{0, 0, 0}{\parbox{\textwidth}{\hspace{15pt}#1#2#3}}}
\captionsetup[lstlisting]{format=listing, labelfont=white, textfont=white, singlelinecheck=false, font={bf,footnotesize}}




\lstset{xleftmargin=\fboxsep, xrightmargin=-\fboxsep}

%\DeclareCaptionFormat{listing}{%
%  \parbox{\textwidth}{\colorbox{gray}{\parbox{\textwidth}{#1#2#3}}\vskip-4pt}}
%\captionsetup[lstlisting]{format=listing,labelfont=white,textfont=white}
%\lstset{frame=lrb,xleftmargin=\fboxsep,xrightmargin=-\fboxsep}
%\lstset{frameround=tttt}


%\renewcommand{\thelstlisting}{\thesection-\arabic{lstlisting}}

\setlength{\marginparwidth}{1.2in}
\let\oldmarginpar\marginpar
\renewcommand\marginpar[1]{\-\oldmarginpar[\raggedleft\footnotesize #1]%
{\raggedright\footnotesize #1}}

\newenvironment{addcode}
{\begin{texttt}{}{\setlength{\leftmargin}{1em}}\item\scriptsize\bfseries}
{\end{texttt}}

%error below


%vers{0.4}{0.5} version between "releases"
%Created by Matz
%Since 31/8 2012, vers clashes with clearpage...
%We cannot compare floats but point metric is fine
\newcommand{\vers}[2]%
 {%
\ifnum 1=0%
     \ifdim #2pt > #1pt%
            \marginpar{\bf New:\\ #1 $\rightarrow$ #2}%
     \else%
            \marginpar{(Ver: #2)}    %
     \fi%
     \ifdim #2pt>1pt %give tips: on or off
            \begin{flushright}%
               % \footnotesize
                \ifdim #2pt > 0.2pt%
                       \colorbox{lightblue}{Check content of text}
                \fi%
                \ifdim #2pt > 0.5pt%
                       \kommentar{Check that citations and structure of the text is correct}%
                \fi%

                \ifdim #2pt > 0.8pt%
                       \kommentar{Check that grammar and punctuation is correct}%
                \fi%

                \ifdim #2pt > 0.94pt%
                       \kommentar{This text is supposed to be the final version}%
                \fi%
            \end{flushright}%
     \fi%
\fi%
}



\newcommand{\about}[1]
 {
\footnote{About: #1}
 }


\newcommand{\ver}[1]
 {
%#1pt
%\vspace{-2Em}
%\marginpar{Version: #1}
\vers{#1}{#1} %no change between versions
 }


%\DeclareBibliographyCategory{Book}
%\DeclareBibliographyCategory{Report}
%\DeclareBibliographyCategory{Web page}
%\DeclareBibliographyCategory{Procedings}

\def \frw {9} %figure result width in centimeter

\def \naive {na\"{\i}ve}

\def \codecushion {\vspace{-2Em}}
\def \tablecapdist {\vspace{2Em}}
\def \tabletopdist {\vspace{1Em}}

\def \eg {e.g.,\ }
\def \ie {i.e.,\ }

%Brittish English use without z
\def \realize {realise }
\def \realized {realised }

\def \clearpages {\clearpage }

%\frenchspacing

\hyphenation{single-Comp-Thread}
%\usepackage{url}%break path
%http://tex.stackexchange.com/questions/299/how-to-get-long-texttt-sections-to-break

\begin{comment}
\newcommand*\justify{%
  \fontdimen2\font=0.4em% interword space
  \fontdimen3\font=0.2em% interword stretch
  \fontdimen4\font=0.1em% interword shrink
  \fontdimen7\font=0.1em% extra space
  \hyphenchar\font=`\-% allowing hyphenation
}
\end{comment}


%error above

%http://www.texample.net/tikz/examples/framed-tikz/

\usepackage{framed}

%http://www.forkosh.com/pstex/latexcommands.htm
%\renewcommand{\rmdefault}{garamond}

%\newcommand{\tmpsection}[1]{}




%\showcitenames
%\nocite{*}

%Fixa
\pdfinfo{
   /Author (Matz Johansson Bergstrom)
   /Title  (Creating a PDF document using PDFLaTeX)
   /Subject (PDFLaTeX)
   /Keywords (Convolution;Fast Fourier Transform; GPU; CPU; Multithreading; CUDA)
}

\definecolor{llightgray}{rgb}{0.9, 0.99, 0.93}
\definecolor{tabletitle}{rgb}{0.85, 0.92, 0.85}
\definecolor{tablegray}{rgb}{0.7, 0.8, 0.7}

\definecolor{tablelight}{rgb}{0.98, 0.99, 0.98} %table background

\definecolor{tabledark}{rgb}{0.95, 0.9, 0.8}%not used?



\definecolor{lightgray}{gray}{0.95}
\usepackage[table]{xcolor} %must be loaded last
%F�r att konvertera figurer som har problem med transparans:

%http://docupub.com/pdfconvert/done.aspx?cid=77689d67-6e7f-4364-97d2-d0f7aa245ef0
%konvertera till pdf 1.3 s� skall det ordna sig
\renewcommand*{\glspostdescription}{}

%\psfrag{text in eps to replace}[l][l]{text in latex to use}


%http://www.bradleymedia.org/latex-formatting/
%\titleformat*{\section}
%   {\normalfont\Large\bfseries\color{black} }



\newenvironment{widepar}%
  {\setlength{\rightskip}{-\marginparsep}\addtolength{\rightskip}{-\marginparwidth}}{\par}

%\usepackage{showkeys} %to see the labels

\pdfminorversion=4


%\usepackage{url}

%% Define a new 'leo' style for the package that will use a smaller font.
\makeatletter
\def\url@leostyle{%
  \@ifundefined{selectfont}{\def\UrlFont{\sf}}{\def\UrlFont{\small\ttfamily}}}
\makeatother
%% Now actually use the newly defined style.
\urlstyle{leo}

%\usepackage[hyphens]{url}

\headheight=14pt %fixing fancyheader warning
%\pretolerance=15000 %fixing overfull warning

%\renewcommand{\cftdot}{}

%http://tex.stackexchange.com/questions/38536/how-can-i-pass-underscore-to-newcommand-properly

\newcommand{\codelst}{\begingroup
  \catcode`_=12 \docodelst}
\newcommand{\docodelst}[1]{%
  \lstinputlisting[caption=\texttt{#1}]{#1}%
  \endgroup
}


\renewcommand{\arraystretch}{1.1}
